\section{Finding Agarttha}

In \textbf{King of the World}, Guenon documents Christianity's loss of its primordial Tradition as the gradual break of the link to its spiritual centre. This break occurred in stages following the period called the Middle Ages. The first break came with the destruction of the Templars, since the Orders of Chivalry maintained a spiritual connection between the East and the West. The Rosicrucians made the effort to continue this liaison. The Renaissance and Reformation completed the break, and the Rosicrucians allegedly retired to Asia. Thus, according to Guenon, no western organisation has maintained full and effective initiatory knowledge. So where does Guenon suggest to look? He turns to two clairvoyants, \textbf{Emmanuel Swedenborg} and \textbf{Blessed Anne-Catherine Emmerich}.

\begin{quotex}
According to Swedenborg the `lost Word' must henceforth be sought among the \textbf{Sages of Tibet} and of \textbf{Tartary}, where the mysterious \textbf{Mount of Prophets} of Anne-Catherine Emmerich's vision is also set.
\end{quotex}


Keep in mind that Guenon is always speaking in terms of knowledge, not its outer form. Don't confuse the seed with the shell.


\flrightit{Posted on 2010-08-25 by Cologero}

